\chapter{Funções lineares e afins}\label{ch:g9:linfunc}

``Três quilogramas custam três vezes o que custa um'': isso é
\omterm{def:g9:linfunc:linear}{proporcionalidade}, o jeito mais simples de duas grandezas se
relacionarem. As \omterm{def:g9:linfunc:linear}{funções lineares} são a \omterm{def:g9:linfunc:linear}{proporcionalidade} na linguagem
das funções; as \omterm{def:g9:linfunc:affine}{funções afins} acrescentam uma quantia inicial fixa. Os
gráficos delas são retas, e ler essas retas é a habilidade que este
capítulo constrói.

\section{Proporcionalidade e funções lineares}

\begin{definition}[Função linear]\label{def:g9:linfunc:linear}
Uma \emph{função linear}\index{função linear} multiplica todo número
por um número fixo $a$:
\[
f(x) = a x .
\]
O número $a$ é o \emph{coeficiente} de $f$. Duas grandezas são
\emph{proporcionais}\index{proporcionalidade} exatamente quando uma é
função linear da outra.
\end{definition}

\begin{example}\label{ex:g9:linfunc:linear}
Se as maçãs custam $3$ por quilograma, o preço de $x$ quilogramas é
$f(x) = 3x$: comprar o dobro custa o dobro, e o preço de $2.5$ kg é
$f(2.5) = 7.5$.

Reciprocamente, se uma \omterm{def:g9:linfunc:linear}{função linear} satisfaz $f(4) = 10$, então
$a = \frac{10}{4} = 2.5$ e $f(x) = 2.5x$: \emph{um único valor
determina uma \omterm{def:g9:linfunc:linear}{função linear} por completo.}
\end{example}

\begin{proposition}[Gráfico de uma função linear]\label{prop:g9:linfunc:lineargraph}
O gráfico de $f(x) = ax$ é uma reta \emph{que passa pela origem}.
Reciprocamente, toda reta não vertical que passa pela origem é o
gráfico de uma \omterm{def:g9:linfunc:linear}{função linear}.
\end{proposition}
\admitted

\begin{omfigure}
\begin{tikzpicture}
\begin{axis}[omaxis,
  width=7.4cm, height=5.6cm,
  xmin=-3.4, xmax=3.4, ymin=-3.2, ymax=3.6,
  xtick={-2,-1,1,2}, ytick={-2,-1,1,2}]
  \addplot[omDef, thick, domain=-2.2:2.2] {1.5*x};
  \addplot[omThm, thick, domain=-3.1:3.1] {-0.5*x};
  \node[omDef, font=\small] at (axis cs:2.15,2.6) {$y = 1.5x$};
  \node[omThm, font=\small] at (axis cs:-2.5,1.7) {$y = -0.5x$};
\end{axis}
\end{tikzpicture}

{\small Duas \omterm{def:g9:linfunc:linear}{funções lineares}: os dois gráficos passam pela origem; o
coeficiente $a$ é a altura alcançada em $x = 1$.}
\end{omfigure}

\section{Funções afins}

\begin{definition}[Função afim]\label{def:g9:linfunc:affine}
Uma \emph{função afim}\index{função afim} tem a forma
\[
f(x) = a x + b ,
\]
em que $a$ e $b$ são números fixos. As \omterm{def:g9:linfunc:linear}{funções lineares} são o caso
particular $b = 0$; as funções constantes, o caso $a = 0$.
\end{definition}

\begin{example}
Uma academia cobra uma taxa de matrícula de $20$ e depois $5$ por
visita: o custo total de $x$ visitas é $f(x) = 5x + 20$. O custo
\emph{não} é proporcional ao número de visitas ($10$ visitas não custam
o dobro de $5$), mas cada visita a mais acrescenta a mesma quantia,
$5$.
\end{example}

\begin{proposition}[Gráfico de uma função afim]\label{prop:g9:linfunc:affinegraph}
O gráfico de $f(x) = ax + b$ é uma reta:
\begin{itemize}
  \item $b = f(0)$ é o \emph{coeficiente linear}: a reta cruza o eixo
        vertical em $(0, b)$;
  \item $a$ é o \emph{coeficiente angular}\index{coeficiente angular}:
        andar uma unidade para a direita sobe $a$ unidades (desce, se
        $a < 0$); a função é crescente quando $a > 0$ e decrescente
        quando $a < 0$;
  \item para quaisquer duas entradas $u \neq v$,
        $a = \dfrac{f(v) - f(u)}{v - u}$.
\end{itemize}
\end{proposition}

\begin{proof}[Demonstração do último ponto]
$f(v) - f(u) = (av + b) - (au + b) = a(v - u)$; divida por $v - u$. O
resto é admitido neste nível (o volume do ensino médio dá um tratamento
completo das retas e das \omterm{def:g8:equations:def}{equações} delas).
\end{proof}

\begin{omfigure}
\begin{tikzpicture}
\begin{axis}[omaxis,
  width=7.8cm, height=5.8cm,
  xmin=-1.9, xmax=4.4, ymin=-1.4, ymax=5.4,
  xtick={1,2,3}, ytick={1,2,3}]
  \addplot[omDef, thick, domain=-1.6:4.1] {x + 1};
  \fill (axis cs:0,1) circle (1.5pt);
  \node[left, font=\small] at (axis cs:0,1.25) {$b$};
  \draw[dashed, omThm] (axis cs:1,2) -- (axis cs:2,2) -- (axis cs:2,3);
  \node[font=\footnotesize, below] at (axis cs:1.5,2) {$+1$};
  \node[font=\footnotesize, right] at (axis cs:2,2.5) {$+a$};
\end{axis}
\end{tikzpicture}

{\small Ler $y = ax + b$ no gráfico: a reta cruza o eixo vertical na
altura $b$ e sobe $a$ a cada passo de uma unidade para a direita (aqui,
$a = 1$, $b = 1$).}
\end{omfigure}

\begin{method}[Achar uma função afim a partir de dois valores]\label{met:g9:linfunc:findaffine}
Dados $f(u)$ e $f(v)$ com $u \neq v$:
\begin{enumerate}
  \item calcule o \omterm{prop:g9:linfunc:affinegraph}{coeficiente angular}
        $a = \dfrac{f(v) - f(u)}{v - u}$;
  \item substitua um dos dois valores conhecidos em $f(x) = ax + b$
        para achar $b$;
  \item confira com o outro valor.
\end{enumerate}
\end{method}

\begin{example}\label{ex:g9:linfunc:findaffine}
Ache a \omterm{def:g9:linfunc:affine}{função afim} com $f(2) = 7$ e $f(5) = 16$.
\[
a = \frac{16 - 7}{5 - 2} = \frac93 = 3 ,
\]
e depois $f(2) = 3 \times 2 + b = 7$ dá $b = 1$: $f(x) = 3x + 1$.
Conferindo: $f(5) = 15 + 1 = 16$.
\end{example}

\section{Porcentagens}

\begin{proposition}[Variação percentual]\label{prop:g9:linfunc:percent}
Aumentar uma grandeza em $t\,\%$ a multiplica por
$1 + \dfrac{t}{100}$; diminuí-la em $t\,\%$ a multiplica por
$1 - \dfrac{t}{100}$. As variações percentuais são \omterm{def:g9:linfunc:linear}{funções lineares}.
\end{proposition}

\begin{proof}
Aumentar $x$ em $t\,\%$ significa somar $\frac{t}{100}\,x$:
\[
x + \frac{t}{100}\,x = \left(1 + \frac{t}{100}\right) x .
\]
Mesma conta, com sinal de menos, para uma diminuição.
\end{proof}

\begin{example}\label{ex:g9:linfunc:percent}
Um preço de $80$ aumenta $15\%$: preço novo
$80 \times 1.15 = 92$. Um preço de $60$ diminui $30\%$:
$60 \times 0.7 = 42$.

\emph{Encadear variações multiplica os fatores.} Um aumento de $20\%$
seguido de uma redução de $20\%$ dá $1.2 \times 0.8 = 0.96$: uma
\emph{queda} de $4\%$ no total, e não uma volta ao ponto de partida!
\end{example}

\section{Exercícios}

\begin{exercise}[$\star$]\label{exo:g9:linfunc:1}
Entre estas funções, quais são lineares? Quais são afins? Quais não são
nem uma coisa nem outra?
\[
f(x) = 4x, \qquad
g(x) = 3x - 5, \qquad
h(x) = x^2, \qquad
k(x) = -\tfrac x2, \qquad
m(x) = 7 .
\]
\end{exercise}

\begin{exercise}[$\star$]\label{exo:g9:linfunc:2}
Seja $f(x) = 2x - 3$. Calcule $f(0)$, $f(4)$, $f(-1)$ e ache o número
$x$ tal que $f(x) = 9$.
\end{exercise}

\begin{exercise}[$\star$]\label{exo:g9:linfunc:3}
Uma \omterm{def:g9:linfunc:linear}{função linear} satisfaz $f(6) = 15$. Ache o coeficiente dela e
depois calcule $f(10)$ e $f(-2)$.
\end{exercise}

\begin{exercise}[$\star$]\label{exo:g9:linfunc:4}
Trace no mesmo sistema de eixos os gráficos de $f(x) = 2x$,
$g(x) = 2x + 3$ e $h(x) = -x + 4$. O que se pode dizer dos gráficos de
$f$ e de $g$?
\end{exercise}

\begin{exercise}[$\star$]\label{exo:g9:linfunc:5}
Ache a \omterm{def:g9:linfunc:affine}{função afim} tal que $f(1) = 5$ e $f(3) = 11$; e depois a função
tal que $g(0) = 4$ e $g(2) = 0$.
\end{exercise}

\begin{exercise}[$\star\star$]\label{exo:g9:linfunc:6}
Calcule: o preço depois de um aumento de $25\%$ sobre $120$; o preço
depois de um desconto de $40\%$ sobre $65$; o preço original se um
artigo custa $69$ depois de um desconto de $25\%$.
\end{exercise}

\begin{exercise}[$\star\star$]\label{exo:g9:linfunc:7}
Um plano de celular A custa $10$ por mês mais $0.05$ por minuto de
ligação; o plano B custa $25$ por mês com ligações ilimitadas.
\begin{enumerate}
  \item Expresse o custo mensal do plano A como função do número $x$ de
        minutos.
  \item A partir de quantos minutos por mês o plano B fica mais barato?
\end{enumerate}
\end{exercise}

\begin{exercise}[$\star\star$]\label{exo:g9:linfunc:8}
Uma população de $2000$ bactérias aumenta $10\%$ a cada hora.
\begin{enumerate}
  \item Quantas bactérias há depois de uma hora? E depois de duas
        horas?
  \item Explique por que a resposta depois de duas horas não é $2400$.
\end{enumerate}
\end{exercise}

\begin{exercise}[$\star\star$]\label{exo:g9:linfunc:9}
O gráfico de uma \omterm{def:g9:linfunc:affine}{função afim} passa pelos pontos $(2, 3)$ e $(6, 1)$.
\begin{enumerate}
  \item Calcule o \omterm{prop:g9:linfunc:affinegraph}{coeficiente angular} dela. A função é crescente ou
        decrescente?
  \item Dê a expressão dela e calcule a entrada para a qual a saída é
        $0$.
\end{enumerate}
\end{exercise}

\begin{exercise}[$\star\star\star$]\label{exo:g9:linfunc:10}
Uma loja primeiro aumenta um preço em $t\,\%$ e depois reduz o preço
novo em $t\,\%$.
\begin{enumerate}
  \item Mostre que o preço final é igual ao original multiplicado por
        $1 - \dfrac{t^2}{10000}$.
  \item Deduza que o preço final é sempre \emph{menor} que o original
        (para $0 < t \leq 100$) e ache a variação percentual total para
        $t = 10$.
\end{enumerate}
\end{exercise}

\section{Problema: os dois termômetros}

\begin{problem}[{Problema de fim de semana --- Celsius, Fahrenheit e as
funções afins do dia a dia: uma temperatura marca o mesmo número nas
duas escalas, e o coeficiente angular constante é a assinatura da
retidão}]
\label{pb:g9:linfunc:1}
Em algumas partes do mundo, ouve-se ``$68$ graus'' e se veste roupa
leve; em outras, ouve-se o mesmo e se pega um casaco. Os dois
termômetros do planeta estão ligados por uma \omterm{def:g9:linfunc:affine}{função afim} --- e
construí-la a partir de dois fatos é exatamente a habilidade do
\cref{met:g9:linfunc:findaffine}. Este problema constrói a conversão,
encontra a temperatura misteriosa em que as duas escalas concordam,
depois lê táxis, grilos e velas com os mesmos óculos e termina com o
teste que reconhece a própria retidão.

\medskip
\noindent\textbf{Parte I --- Construir a conversão.}
Dois fatos fixam a escala Fahrenheit: a água congela a
$32\,^\circ$F e ferve a $212\,^\circ$F (a $0\,^\circ$C e
$100\,^\circ$C).
\begin{enumerate}
  \item Ache a \omterm{def:g9:linfunc:affine}{função afim} $F = f(C)$ que converte Celsius em
        Fahrenheit (\cref{met:g9:linfunc:findaffine} com os valores
        $f(0) = 32$ e $f(100) = 212$).
  \item Converta para Fahrenheit: um dia ameno de $20\,^\circ$C; os
        $37\,^\circ$C do corpo; um gélido $-10\,^\circ$C.
  \item Isole $C$ para construir a conversão de volta e use-a em
        $50\,^\circ$F e em $98.6\,^\circ$F.
  \item $f$ é uma função \emph{linear}
        (\cref{def:g9:linfunc:linear})? Teste os sinais reveladores:
        quanto vale $f(0)$, e dobrar a leitura em Celsius dobra a em
        Fahrenheit? Conclua com a palavra certa
        (\cref{def:g9:linfunc:affine}).
  \item A adivinhação clássica: existe uma temperatura que marca
        \emph{o mesmo número} nas duas escalas? Resolva $f(x) = x$ e
        descreva o que a sua solução significa no gráfico de $f$ (que
        reta o gráfico cruza ali?).
\end{enumerate}

\medskip
\noindent\textbf{Parte II --- \omterm{def:g9:linfunc:affine}{Funções afins} soltas na natureza.}
\begin{enumerate}[resume]
  \item O táxi A cobra $3$ euros mais $1.50$ por km; o táxi B cobra
        $2$ euros por km, sem taxa fixa. Escreva as duas funções de
        preço e compare-as para uma corrida de $4$ km e uma de $8$ km.
  \item Ache a distância de equilíbrio em que os dois táxis custam o
        mesmo e descreva a situação num gráfico (duas retas --- o que
        acontece no ponto de cruzamento, antes dele e depois dele?).
  \item Regra prática dos naturalistas: um grilo cricrila cerca de
        $N = 7T - 30$ vezes por minuto à temperatura $T\,^\circ$C.
        Quantos cricrilos numa noite de $20\,^\circ$C? Você conta $82$
        cricrilos num minuto: que temperatura o grilo anuncia?
  \item Um celular comprado por $600$ euros perde $15$ euros de valor
        por mês: escreva a função de valor $v(t)$, ache quando o
        celular vale $240$ euros e quando o modelo diz que ele não vale
        mais nada. A fórmula ainda significa alguma coisa depois dessa
        data?
  \item As variações percentuais são \omterm{def:g9:linfunc:linear}{funções lineares}: $+25\,\%$ é
        multiplicar por $1.25$, e $-20\,\%$, por $0.8$
        (\cref{prop:g9:linfunc:percent}). Componha as duas: o que um
        aumento de $+25\,\%$ seguido de um corte de $-20\,\%$ faz com
        um preço? Explique o milagrezinho e compare com o
        \cref{exo:g9:linfunc:10}.
\end{enumerate}

\medskip
\noindent\textbf{Parte III --- O \omterm{prop:g9:linfunc:affinegraph}{coeficiente angular} é tudo.}
\begin{enumerate}[resume]
  \item Uma vela mede $20$ cm quando acesa e $17$ cm meia hora depois.
        Supondo derretimento afim, ache $h(t)$ (altura em cm, $t$ em
        minutos) e preveja quando a vela se apaga.
  \item Dois ciclistas percorrem a mesma estrada: a distância do
        primeiro é $d_1(t) = 24t$ (km depois de $t$ horas), a do
        segundo é $d_2(t) = 10 + 18t$. Quem partiu com vantagem, quem
        pedala mais rápido e em que instante e em que quilômetro o
        primeiro alcança o segundo?
  \item Sejam $f(x) = 2x + 3$ e $g(x) = -0.5x + 7$. Calcule
        $f(x) + g(x)$ e verifique que ela é afim de novo. Em geral,
        quais são o \omterm{prop:g9:linfunc:affinegraph}{coeficiente angular} e o coeficiente linear de uma
        \omterm{def:g1:addition:def}{soma}?
  \item Discuta completamente a \omterm{def:g8:equations:def}{equação} $mx + p = 0$: quantas soluções
        quando $m \neq 0$ (e qual)? Quando $m = 0$ e $p \neq 0$? Quando
        $m = 0$ e $p = 0$? (Cada caso tem uma história gráfica: onde uma
        reta de \omterm{prop:g9:linfunc:affinegraph}{coeficiente angular} $m$ cruza o eixo horizontal?)
  \item O teste do alinhamento: os pontos $(1, 5)$, $(3, 9)$,
        $(7, 17)$ estão sobre uma mesma reta? E $(1, 5)$, $(3, 9)$,
        $(6, 16)$? Calcule as taxas de variação entre os pares e
        decida. Enuncie a moral: uma \emph{taxa de variação constante}
        é a assinatura das \omterm{def:g9:linfunc:affine}{funções afins} --- os gráficos curvos mudam
        de taxa, e medir \emph{isso} é assunto do volume do ensino
        médio.
\end{enumerate}
\end{problem}
